% ===== 原文 pp.33–46（PDF p.48–61）=====
% 译者：tr-02（第 2 章 §2.1 前段；tag 区间 2.1–2.7；本章标题由本块书写）
% [原文 p.33]
\section{A-模型与B-模型}

本章的目标是说明镜像对称故事中所涉及的 A-模型与 B-模型，以满足本书之需。这里选取的题目颇具选择性。A-模型关心的是 Gromov–Witten 不变量。关于 Gromov–Witten 不变量的定义，我们将极为简略——其他文献已对此作了详尽考察。另一方面，我们将详细讨论 Gromov–Witten 不变量所产生的某些结构，例如半无限 Hodge 结构变分这一概念。B-模型关心的是周期积分。对 B-模型，我们将完全略过 Calabi–Yau 情形，读者可在许多文献（包括 [43]）中找到相关讨论。我们转而关注本书主要感兴趣的情形，即 Fano 流形的情形。同样，在此情形下，相关的结构也是半无限 Hodge 结构变分。

% [Tutor Note: 原文 §2.1 之下还有 2.1.1/2.1.2/2.1.3 三级标题，本译用 \subsubsection 呈现（自动编号恰为 2.1.1 等），与原文编号一致；若全稿统一要求用 \para 可回退。另：原文 "variations of Hodge structure" 全书沿用 "Hodge 结构变分"，与 glossary 的 "Hodge 结构" 配套。]

\subsection{A-模型}

\subsubsection{稳定映射与 Gromov–Witten 不变量}

我们从基本定义开始。

\begin{definition}
一个\emph{稳定 $n$ 点曲线}是指数据 $(C,x_1,\dots,x_n)$，其中 $C$ 是代数闭域上的一条（可能可约的）真、既约、连通代数曲线，其奇点仅为节点，而 $x_1,\dots,x_n\in C$ 是互不相同且不与任何节点重合的点。此外，$(C,x_1,\dots,x_n)$ 的自同构群必须有限。这里，$(C,x_1,\dots,x_n)$ 的自同构是指满足 $\varphi(x_i)=x_i$ 的自同构 $\varphi:C\to C$。$C$ 的\emph{亏格}是 $C$ 的算术亏格。$n$ 点曲线 $(C,x_1,\dots,x_n)$ 与 $(C',x_1',\dots,x_n')$ 之间的同构是指满足 $\varphi(x_i)=x_i'$ 的同构 $\varphi:C\to C'$。
\end{definition}

我们指出，$(C,x_1,\dots,x_n)$ 具有有限自同构群这一条件可以刻画如下。设 $\nu:\tilde C\to C$ 为 $C$ 的正规化，称 $\tilde C$ 中的一点为\emph{区分点}（distinguished point），若它经 $\nu$ 或者映到某个节点、或者映到某个标记点 $x_1,\dots,x_n$。于是 $(C,x_1,\dots,x_n)$ 具有有限自同构群，等价于：$\tilde C$ 的每个亏格零分支至少有 3 个区分点，且 $\tilde C$ 的每个亏格一分支至少有 1 个区分点。

Deligne 与 Mumford [19] 的一个著名结果是：$n$ 点稳定曲线的模空间 $\overline{\cM}_{g,n}$ 是不可约、真、光滑的 Deligne–Mumford 叠。这句话需要一些讨论。我们将采取 [56] 第 139–149 页所展示的对待叠的态度。换言之，我们不给出叠的定义，而只是尽量不过分担心这一缺口。然而要点如下。在一个理想世界里，$\overline{\cM}_{g,n}$ 会是一个精细模空间，即会是一个概形，表示从概形范畴到
% [原文 p.34]
集合范畴的以下函子
\[
S\mapsto\left\{
\begin{array}{l}
\text{平坦族 $\cC\to S$ 的同构类，带有截面 $\sigma_1,\dots,\sigma_n:S\to\cC$，}\\
\text{使得对 $S$ 的每个几何点 $\bar s$，$(\cC_{\bar s},\sigma_1(\bar s),\dots,\sigma_n(\bar s))$}\\
\text{都是稳定 $n$ 点亏格 $g$ 曲线}
\end{array}
\right\}.
\]
于是 $\operatorname{Hom}(S,\overline{\cM}_{g,n})$ 会与 $S$ 上这类平坦族的同构类集合一致。特别地，$\overline{\cM}_{g,n}$ 上会存在稳定 $n$ 点曲线的万有族，它对应于 $\operatorname{Hom}(\overline{\cM}_{g,n},\overline{\cM}_{g,n})$ 中的恒等元。然而众所周知，由于曲线自同构的存在，并不存在这样的万有族：稳定曲线的自同构群有限，但未必平凡。于是必须扩大概形范畴，使模函子变得可表示。这通过使用代数叠范畴来完成。光滑 Deligne–Mumford 叠可以看作一个代数几何对象，它是轨道流形，即一个概形，局部是一个光滑概形被一个有限群取商，但我们记住这一局部描述。我们仍然可以谈论 Chow 群（或若在 $\C$ 上按通常拓扑工作，则谈论上同调群）与相交积之类的东西，但它们至多在 $\Q$ 上有定义。

% [Tutor Note: 原文 "as best defined over Q" 指 Deligne–Mumford 叠上 Chow 群只给出 $\Q$ 系数（原文前一句正是 "these are at best defined over Q"），故下一段才强调 "intersection numbers are only defined over Q"；这不是笔误，勿改。]

特别值得强调的是：相交数只在 $\Q$ 上有定义。事实上，若局部可以把光滑 Deligne–Mumford 叠描述为 $X/G$，其中 $X$ 光滑，那么非形式地说，若 $Z_1,Z_2\subseteq X/G$ 是两个补维数的子概形、相交于一个零维子概形，则 $Z_1$ 与 $Z_2$ 在 $P\in Z_1\cap Z_2$ 处的相交数定义为 $\frac{1}{|G|}\sum_{Q\in\pi^{-1}(P)}i_Q$，其中 $\pi:X\to X/G$ 是投影，$i_Q$ 表示 $\pi^{-1}(Z_1)$ 与 $\pi^{-1}(Z_2)$ 在 $Q$ 处的相交重数。

例如，取 $X=\PP^2$，$G=\Z_2$，$\Z_2$ 通过
\[
(x,y,z)\mapsto(-x,-y,z)
\]
作用。我们可以把 $X/G$ 看作一个光滑 Deligne–Mumford 叠。设 $D_1$ 与 $D_2$ 分别是 $X$ 中由 $x=0$ 与 $y=0$ 给出的直线 $L_1$ 与 $L_2$ 在 $X/G$ 中的像。若 $\pi:X\to X/G$ 表示投影，则我们有
\[
D_1\cdot D_2=\frac{1}{2}L_1\cdot L_2=\frac{1}{2}.
\]

由此，给定一个光滑 Deligne–Mumford 叠 $X$，我们就得到 $A^i_{\Q}(X)$，即带有理系数的余维 $i$ Chow 群，以及相交积
\[
A^i_{\Q}(X)\times A^j_{\Q}(X)\to A^{i+j}_{\Q}(X),
\]
它给出 $A^{*}_{\Q}(X)$ 上的环结构。类似地，若在 $\C$ 上按通常拓扑工作，则有意义的是有理上同调环 $H^{*}(X,\Q)$，以及杯积
\[
H^i(X,\Q)\times H^j(X,\Q)\to H^{i+j}(X,\Q).
\]
我们通常在后一种框架下工作。

\begin{examples}
当 $n\le 2$ 时 $\overline{\cM}_{0,n}$ 为空，而 $n=3$ 时它只是一个点。事实上，$\overline{\cM}_{0,n}$ 总是概形，因为 $n\ge 3$ 的 $n$ 点有理曲线没有自同构。$\overline{\cM}_{0,4}$ 同构于 $\PP^1$，这可见如下。给定 $(\PP^1,x_1,x_2,x_3,x_4)$，存在唯一的 $\operatorname{PGL}(2)$ 元素把 $x_1$ 送到 $0$、$x_2$ 送到 $1$、$x_3$ 送到 $\infty$。于是
% [原文 p.35]
$x_4\in\PP^1\setminus\{0,1,\infty\}$ 决定了 $(\PP^1,x_1,x_2,x_3,x_4)$ 的同构类。（这不过是经典的交比。）因此，$\overline{\cM}_{0,4}$ 中对应于非奇异曲线的开子集 $\cM_{0,4}$ 同构于 $\PP^1\setminus\{0,1,\infty\}$。点 $0,1,\infty$ 对应于图 1 所绘的稳定但奇异的曲线。
\end{examples}

\begin{figure}[htbp]
\centering
\includegraphics[width=0.85\textwidth]{f2_1.png}
\figtitle{1}{三条稳定四点曲线的构型}
\end{figure}

% [Tutor Note: 原文 Figure 1 只有图题 "FIGURE 1"，无说明文字；中译图注仅为一句简短图意说明，不添新内容。三条曲线分别对应交比参数 $0,1,\infty$：原文在正文中点明 "The points 0,1,∞ correspond to stable but singular curves depicted in Figure 1"，可作核对锚点；图上标签 $x_1,\dots,x_4$ 与 $0,1,\infty$ 保持原图。]

注意 $\overline{\cM}_{0,4}=\PP^1$ 上的万有族可这样得到：取由投影到第二个分量给出的族 $\PP^1\times\PP^1\to\PP^1$，并把 $\{0\}\times\PP^1$、$\{1\}\times\PP^1$ 与 $\{\infty\}\times\PP^1$ 以及对角截面分别取为 $\sigma_1,\dots,\sigma_4$。然后我们爆破 $\sigma_4$ 与另外三个截面之一相交的那三个点。

按传统的看法，这说的是：当 $x_4$ 逼近另外三个标记点之一时，我们必须吹出一个 $\PP^1$ 副本，以保持标记点互不相同。

$\overline{\cM}_{1,1}$ 可以等同于 $j$-直线，但不是作为概形。每条带标记椭圆曲线都有一个由取负给出的自同构，而 $j$-不变量为 $0$ 与 $1$ 的椭圆曲线有更大的自同构群。因此，即使是 $j=0$ 与 $j=1$ 之外的点，也必须看作叠点，局部形如 $U/\Z_2$，其中 $\Z_2$ 在 $U$ 上平凡地作用！

一般地，$\dim\overline{\cM}_{g,n}=3g-3+n$。

下面考虑稳定映射。

\begin{definition}
设 $X$ 是一个簇。一个到 $X$ 的\emph{稳定 $n$ 点映射}是一个映射
\[
f:(C,x_1,\dots,x_n)\to X,
\]
其中 $C$ 是真、连通、既约的节点代数曲线，$x_1,\dots,x_n\in C$ 是互不相同且不与 $C$ 的任何节点重合的点，并且 $f$ 的自同构群有限。这里，$f$ 的自同构是指 $(C,x_1,\dots,x_n)$ 的自同构 $\varphi$，满足 $f\circ\varphi=f$。
\end{definition}

注意，$f$ 的自同构群有限，等价于：$C$ 的正规化的每个使 $f$ 为常值的分支，若亏格为零则至少有 3 个区分点，若亏格为一则至少有 1 个区分点。

若 $\beta\in H_2(X,\Z)$，我们说 $f$ \emph{表示} $\beta$，是指 $f_*([C])=\beta$，其中 $[C]\in H_2(C,\Z)$ 是基本类。（若不在域 $\C$ 上工作，可以把 $H_2$ 换成维数为一的 Chow 群。）

% [原文 p.36]
两个稳定映射 $f:(C,x_1,\dots,x_n)\to X$、$f':(C',x_1',\dots,x_n')\to X$ 视为同构，若存在同构 $\varphi:(C,x_1,\dots,x_n)\to(C',x_1',\dots,x_n')$ 使得 $\varphi(x_i)=x_i'$ 且 $f'\circ\varphi=f$。

与前面一样，我们可以问模函子
\[
S\mapsto\left\{
\begin{array}{l}
\text{平坦族 $\cC\to S$ 的同构类，带有截面 $\sigma_1,\dots,\sigma_n:S\to\cC$ 与态射 $f:\cC\to X$，}\\
\text{使得对 $S$ 的每个几何点 $\bar s$，$f:(\cC_{\bar s},\sigma_1(\bar s),\dots,\sigma_n(\bar s))\to X$}\\
\text{都是表示 $\beta$ 的亏格 $g$ 稳定映射}
\end{array}
\right\}
\]
是否可表示。它一般不能由概形表示，甚至不能由光滑 Deligne–Mumford 叠表示。然而它是真 Deligne–Mumford 叠 [9]。这会使我们的处境更为困难。我们把亏格 $g$、表示类 $\beta\in H_2(X,\Z)$ 的 $n$ 点稳定映射的叠记作 $\overline{\cM}_{g,n}(X,\beta)$。

这里还要说明：真性可以用赋值判据来检验，而该判据归结为以下基本事实，称为稳定约化。证明见 [28]，§4.2。

\begin{proposition}
设 $S$ 是 $k$ 上的非奇异曲线，$s\in S$ 是一点。令 $U=S\setminus\{s\}$。假设在 $U$ 上有一个稳定映射的族 $f:(C_U,\sigma_1,\dots,\sigma_n)\to X$。则存在 $s$ 的开邻域 $V$、一个有限映射 $\pi:V'\to V$（其中 $V'$ 是非奇异曲线）以及一点 $s'\in V'$，使得
\begin{enumerate}[label=(\arabic*)]
\item 若 $U'=V'\setminus\{s'\}$，则 $\pi|_{U'}$ 是 étale 的；
\item 令 $C_{U'}:=C_U\times_U U'$。这给出一个拉回的稳定映射族 $f':(C_{U'},\sigma_1',\dots,\sigma_n')\to X$，它由复合 $C_{U'}\to C_U\xrightarrow{f}X$ 给出。那么 $f'$ 可以延拓为 $V'$ 上的稳定映射族 $(C_{V'},\sigma_1',\dots,\sigma_n')\to X$。
\end{enumerate}
\end{proposition}

\begin{examples}
考虑目标空间 $X=\PP^2$，$[\ell]\in H_2(X,\Z)$ 是一条直线的同调类。那么 $\overline{\cM}_{0,0}(X,[\ell])$ 当然就是 $(\PP^2)^{*}$，即 $\PP^2$ 的对偶。另一方面，$\overline{\cM}_{0,1}(X,[\ell])$ 是关联对应 $I\subseteq\PP^2\times(\PP^2)^{*}$，其中
\[
I=\{(x,\ell)\mid x\in\ell\}.
\]
$\overline{\cM}_{0,0}(X,2[\ell])$ 稍微复杂一些。这个模空间中有四类稳定映射。第一，$f$ 的定义域 $C$ 可能不可约，而 $f(C)$ 是一条二次曲线；或者 $C$ 可能是两条直线之并，而 $f(C)$ 是可约二次曲线。其次，$f$ 可以是直线 $\ell\subseteq\PP^2$ 的二重覆盖，此时定义域或者不可约、或者可约。在这两种二重覆盖情形，$f$ 有一个 2 阶非平凡自同构，所以这些点是 $\overline{\cM}_{0,0}(\PP^2,2[\ell])$ 中的叠点。

再考虑 $\overline{\cM}_{1,0}(\PP^2,[\ell])$。不存在从椭圆曲线到 $\PP^2$ 表示类 $[\ell]$ 的映射，所以乍看之下这个模空间似乎是空的。然而，我们有定义域为 $C=C_1\cup C_2$ 的映射 $f$，其中 $C_1$ 是一条 $\PP^1$，$C_2$ 或是一条椭圆曲线、或是一条在一点处附接到 $C_1$ 的节点有理曲线。此外 $f$ 把 $C_1$ 嵌入为一条直线，并在 $C_2$ 上为常值。由于这既需要选取 $C_1$ 在 $\PP^2$ 中的像、又需要选取 $C_1$ 与 $C_2$ 上使二者相交的点，我们看到
\begin{equation*}
\overline{\cM}_{1,0}(\PP^2,[\ell])=\overline{\cM}_{0,1}(\PP^2,[\ell])\times\overline{\cM}_{1,1}. \tag*{$\square$}
\end{equation*}
\end{examples}

这里我们不打算深入细节，但形变理论的一个标准推论告诉我们，这些模空间有一个期望维数。这个期望
% [原文 p.37]
维数的来由是：给定定义域为亏格 $g$ 的 $n$ 点曲线的稳定映射 $f$，可以计算两个空间 $T_f^1$ 与 $T_f^2$，其中 $T_f^1$ 是 $[f]\in\overline{\cM}_{g,n}(X,\beta)$ 处的切空间，$T_f^2$ 是障碍空间——粗略地说，这意味着 $\overline{\cM}_{g,n}(X,\beta)$ 局部由 $T_f^1$ 中的 $\dim T_f^2$ 个方程切出。这意味着 $[f]$ 在 $\overline{\cM}_{g,n}(X,\beta)$ 中的小开邻域的维数介于
\[
\dim T_f^1-\dim T_f^2\quad\text{与}\quad\dim T_f^1
\]
之间。

数 $\dim T_f^1-\dim T_f^2$ 就是 $\overline{\cM}_{g,n}(X,\beta)$ 的\emph{期望}维数，也称为\emph{虚拟}维数；它可以用 Riemann–Roch 计算，得到公式
\[
\dim T_f^1-\dim T_f^2=n+(\dim_{\C}X-3)(1-g)+\int_{\beta}c_1(\cT_X).
\]

例如，$\overline{\cM}_{0,0}(\PP^2,d[\ell])$ 的虚拟维数是 $3d-1$，因为 $c_1(\cT_{\PP^2})=3[\ell]$。这确实是正确维数，并且在此情形下 $\overline{\cM}_{0,0}(\PP^2,d[\ell])$ 是光滑 Deligne–Mumford 叠。另一方面，$\overline{\cM}_{1,0}(\PP^2,[\ell])$ 的期望维数是 3，而我们看到实际上 $\dim\overline{\cM}_{1,0}(\PP^2,[\ell])=4$。所以事实上，这些模空间未必具有期望维数。更一般地，它们未必光滑。

这对我们想做的事情构成了严重问题。例如，若目标是在 Calabi–Yau 三维流形中（于是 $c_1(\cT_X)=0$）数出表示类 $\beta$ 的亏格 $g$ 的曲线条数，则期望维数为零，于是若模空间 $\overline{\cM}_{g,0}(X,\beta)$ 是维数为零的非奇异空间，我们只需数出这个模空间中点的个数即可。然而我们不能指望这总能发生。

历史上，解决这个问题有两种途径。一种是流形切空间 $X$ 上的复结构形变到一个一般的近复结构，对这种结构，这些模空间往往性态更好。这一途径由 Ruan [100] 与 Ruan–Tian [101] 开创。然而这一途径把我们带进辛几何的领域。另一条途径是构造\emph{虚拟基本类}，由 Behrend–Fantechi [8] 与 Li–Tian [74] 开创。由于这里的目标是综述结果，我们只给出构造的结果，即所谓的虚拟基本类
\[
[\overline{\cM}_{g,n}(X,\beta)]^{\mathrm{vir}}\in A_d(\overline{\cM}_{g,n}(X,\beta))\otimes\Q
\]
（或属于 $H_{2d}(\overline{\cM}_{g,n}(X),\Q)$），其中 $d$ 是 $\overline{\cM}_{g,n}(X,\beta)$ 的虚拟维数。

我们可以用这个类来定义 Gromov–Witten 不变量。有自然的映射
\[
\mathrm{ev}_i:\overline{\cM}_{g,n}(X,\beta)\to X,
\]
其中对稳定映射 $f:(C,x_1,\dots,x_n)\to X$，$\mathrm{ev}_i([f])=f(x_i)$。把它们放在一起得到
\[
\mathrm{ev}=\mathrm{ev}_1\times\cdots\times\mathrm{ev}_n:\overline{\cM}_{g,n}(X,\beta)\to X^n.
\]
对类 $\alpha_1,\dots,\alpha_n\in H^{*}(X,\Q)$、$\beta\in H_2(X,\Z)$，我们定义 Gromov–Witten 不变量
\[
\langle\alpha_1,\dots,\alpha_n\rangle_{g,\beta}=\int_{[\overline{\cM}_{g,n}(X,\beta)]^{\mathrm{vir}}}\mathrm{ev}^{*}(\alpha_1\times\cdots\times\alpha_n)\in\Q.
\]
这里所谓的积分，是指把一个上同调类与一个同调类配对。若愿意，这也可以用 Chow 群来做。注意，若
% [原文 p.38]
$\alpha_i\in H^{d_i}(X,\Q)$（我们把这一点写作 $\deg\alpha_i=d_i$），则该积分等于零，除非 $\overline{\cM}_{g,n}(X,\beta)$ 的虚拟维数的两倍等于 $\sum_{i=1}^{n}d_i$，即
\begin{equation}
2\left(n+(\dim_{\C}X-3)(1-g)+\int_{\beta}c_1(\cT_X)\right)=\sum_{i=1}^{n}d_i. \tag{2.1}
\end{equation}

% [Tutor Note: (2.1) 是维数约束；核对锚点：$X=\PP^2$、$\beta=d[\ell]$、$g=0$ 时左边为 $2(n-1+3d)$，而 $\alpha_i=[\mathrm{pt}]\in H^4$ 给出右边 $4n$，解得 $n=3d-1$——正是 p.38 下半页 Example 2.6 中"过 $n$ 个一般点的 $d$ 次有理曲线"应有 $n=3d-1$ 的来历。]

直观上，Gromov–Witten 不变量可以这样理解。固定 $X$ 的子流形 $Y_1,\dots,Y_n$，其 Poincaré 对偶上同调类为 $\alpha_1,\dots,\alpha_n$；那么这个 Gromov–Witten 不变量应理解为 $X$ 中同调类为 $\beta$、亏格为 $g$、且穿过 $Y_1,\dots,Y_n$ 的曲线条数。这受制于上面给出的解释：所谓曲线是指稳定映射；而就“数”而言，我们需要用虚拟基本类来得到实际的数。最后，一般而言这些数只是有理数，也不必为正。

有一类 Gromov–Witten 不变量我们下文需要用到，即所谓的\emph{引力后代不变量}。在 $\overline{\cM}_{g,n}(X,\beta)$ 上有自然的线丛 $\cL_i$，$i=1,\dots,n$，它在点 $[(C,x_1,\dots,x_n)]$ 处的纤维是余切线 $\mathfrak{m}_{x_i}/\mathfrak{m}_{x_i}^2$，其中 $\mathfrak{m}_{x_i}\subseteq\cO_{C,x_i}$ 是极大理想。令
\[
\psi_i=c_1(\cL_i)\in H^2(\overline{\cM}_{g,n}(X,\beta),\Q).
\]
于是对类 $\alpha_1,\dots,\alpha_n\in H^{*}(X,\Q)$、$\beta\in H_2(X,\Z)$，后代 Gromov–Witten 不变量定义为
\[
\langle\psi^{p_1}\alpha_1,\dots,\psi^{p_n}\alpha_n\rangle_{g,\beta}=\int_{[\overline{\cM}_{g,n}(X,\beta)]^{\mathrm{vir}}}\psi_1^{p_1}\cup\cdots\cup\psi_n^{p_n}\cup\mathrm{ev}^{*}(\alpha_1\times\cdots\times\alpha_n)\in\Q.
\]
注意，若 $\deg\alpha_i=d_i$，则该式为零，除非 $\overline{\cM}_{g,n}(X,\beta)$ 的虚拟维数的两倍等于 $\sum_{i=1}^{n}d_i+\sum_{i=1}^{n}2p_i$，即
\begin{equation}
2\left(n+(\dim_{\C}X-3)(1-g)+\int_{\beta}c_1(\cT_X)\right)=\sum_{i=1}^{n}d_i+\sum_{i=1}^{n}2p_i. \tag{2.2}
\end{equation}

\begin{example}
本书主要关注的例子是 $X=\PP^2$、$g=0$；此时对 $d\ge 0$（且当 $d=0$ 时 $n\ge 3$），$\overline{\cM}_{0,n}(\PP^2,d[\ell])$ 是具期望维数的光滑 Deligne–Mumford 叠。事实上可以证明：若取 $\alpha_1,\dots,\alpha_n\in H^4(X,\Q)$ 为一点的 Poincaré 对偶类 $[\mathrm{pt}]$，则当这个数预期有限时，$\langle\alpha_1,\dots,\alpha_n\rangle_{0,d[\ell]}$ 恰好等于过 $\PP^2$ 中一般点 $P_1,\dots,P_n$ 的 $d$ 次有理曲线条数；而用式 (2.1) 可知，预期有限的情形正是 $n-1+3d=2n$，即 $n=3d-1$。于是例如
\[
\langle[\mathrm{pt}],[\mathrm{pt}]\rangle_{0,[\ell]}=1,
\]
因为过两点有唯一一条直线；而
\[
\langle[\mathrm{pt}]^5\rangle_{0,2[\ell]}:=\langle[\mathrm{pt}],[\mathrm{pt}],[\mathrm{pt}],[\mathrm{pt}],[\mathrm{pt}]\rangle_{0,2[\ell]}=1,
\]
因为过 $\PP^2$ 中 5 个点同样有唯一一条二次曲线。再其次
\[
\langle[\mathrm{pt}]^8\rangle_{0,3[\ell]}=12,
\]
这用到了众所周知的事实：$\PP^2$ 中 8 个一般点包含在三次曲线的一个线性系中，而一般的平面三次曲线线性系含有 12 条奇异、从而是有理的三次曲线。
\end{example}

% [原文 p.39]
在此情形下，我们总得到整数；这里没有叠现象。要点实质上在于：此时只有涉及多重覆盖的曲线才对叠现象有贡献，而这类曲线不会穿过一般位置的点集。然而对后代不变量来说情况并非如此；例如可以证明
\[
\langle\psi^4[\mathrm{pt}]\rangle_{0,2[\ell]}=\frac{1}{8}.
\]
（见例 5.17。）

% [Tutor Note: 原文 p.39 的 $\langle\psi^4[\mathrm{pt}]\rangle_{0,2[\ell]}=1/8$ 是后代不变量"非整数"的典型例子，与上一页非后代情形"总得整数"对照。可用 (2.2) 核对：取 $X=\PP^2$、$d=2$、$n=1$、$d_1=4$、$p_1=4$，左边 $2(1+(2-3)(1-0)+2\cdot3)=2(1-1+6)=12$，右边 $4+2\cdot4=12$，一致。]

这里我们不加证明地列出 Gromov–Witten 不变量与后代 Gromov–Witten 不变量最重要的若干性质。详见例如 [28] 与 [18]。

对 Gromov–Witten 不变量，我们有：

\para{基本类公理。}若 $n+2g\ge 4$，或 $\beta\neq 0$ 且 $n\ge 1$，并且 $[X]\in H^0(X,\Q)$ 是 $X$ 的基本类，则
\[
\langle\alpha_1,\dots,\alpha_{n-1},[X]\rangle_{g,\beta}=0.
\]

\para{除子公理。}若 $n+2g\ge 4$，或 $\beta\neq 0$ 且 $n\ge 1$，并且 $\alpha_n\in H^2(X,\Q)$，则
\[
\langle\alpha_1,\dots,\alpha_n\rangle_{g,\beta}=\left(\int_{\beta}\alpha_n\right)\langle\alpha_1,\dots,\alpha_{n-1}\rangle_{g,\beta}.
\]

\para{点映射公理。}对 $g=0$、$\beta=0$，
\[
\langle\alpha_1,\dots,\alpha_n\rangle_{0,0}=
\begin{cases}
\displaystyle\int_X\alpha_1\cup\alpha_2\cup\alpha_3 & \text{若 } n=3,\\[4pt]
0 & \text{其他情形}.
\end{cases}
\]

后代 Gromov–Witten 不变量满足这些公理的推广，以及一些新的公理。特别地，我们有：

\para{基本类公理。}若 $n+2g\ge 4$，或 $\beta\neq 0$ 且 $n\ge 1$，则
\[
\langle\psi^{p_1}\alpha_1,\dots,\psi^{p_{n-1}}\alpha_{n-1},[X]\rangle_{g,\beta}=\sum_{i=1}^{n-1}\langle\psi^{p_1}\alpha_1,\dots,\psi^{p_i-1}\alpha_i,\dots,\psi^{p_{n-1}}\alpha_{n-1}\rangle_{g,\beta},
\]
其中若 $\psi$ 以负幂出现，则右边的相应不变量取为零。

\para{除子公理。}若 $n+2g\ge 4$，或 $\beta\neq 0$ 且 $n\ge 1$，并且 $\alpha_n\in H^2(X,\Q)$，则
\begin{align*}
&\langle\psi^{p_1}\alpha_1,\dots,\psi^{p_{n-1}}\alpha_{n-1},\alpha_n\rangle_{g,\beta}\\
&\quad=\left(\int_{\beta}\alpha_n\right)\langle\psi^{p_1}\alpha_1,\dots,\psi^{p_{n-1}}\alpha_{n-1}\rangle_{g,\beta}\\
&\qquad+\sum_{i=1}^{n-1}\langle\psi^{p_1}\alpha_1,\dots,\psi^{p_i-1}(\alpha_n\cup\alpha_i),\dots,\psi^{p_{n-1}}\alpha_{n-1}\rangle_{g,\beta},
\end{align*}
若某个 $\psi$ 的幂为负，同上附加说明。

\para{点映射公理。}若 $n\le 3$，则
\[
\langle\psi^{\nu_1}\alpha_1,\dots,\psi^{\nu_n}\alpha_n\rangle_{0,0}=0,
\]
除非 $n=3$ 且 $\nu_1=\cdots=\nu_n=0$。

\para{Dilaton 公理。}
\[
\langle\psi[X],\psi^{p_1}\alpha_1,\dots,\psi^{p_n}\alpha_n\rangle_{g,\beta}=(2g-2+n)\langle\psi^{p_1}\alpha_1,\dots,\psi^{p_n}\alpha_n\rangle_{g,\beta}.
\]

% [原文 p.40]
\subsubsection{量子上同调}

来自弦论的一个重要洞见（见 [112]、[110]）是：Gromov–Witten 不变量可用来扰动上同调中通常的杯积。这个新乘积称为\emph{量子乘积}；它保持结合性这一事实随后给出了对 Gromov–Witten 不变量的强限制，特别地完全确定了 $\PP^2$ 的 Gromov–Witten 不变量。

我们先如下定义 $X$ 的\emph{Gromov–Witten 势}。选取 $H^{*}(X,\C)$ 的一组基 $T_0,\dots,T_m$。我们将始终采用这样的约定：$T_0\in H^0(X,\C)$ 是 $X$ 的基本类，而 $T_1,\dots,T_p$ 生成 $H^2(X,\C)$。令 $\gamma=\sum_{i=0}^{m}y_iT_i$。于是我们定义 Gromov–Witten 势为
\[
\Phi=\sum_{n=0}^{\infty}\sum_{\beta\in H_2(X,\Z)}\frac{1}{n!}\langle\gamma^n\rangle_{0,\beta}.
\]
这里 $\langle\gamma^n\rangle_{0,\beta}=\langle\gamma,\dots,\gamma\rangle_{0,\beta}$，其中 $\gamma$ 取 $n$ 次。于是对给定的 $n$ 与 $\beta$，我们得到一项，它是 $y_0,\dots,y_m$ 的 $n$ 次单项式之和。若对给定的 $n$ 只有有限多个 $\beta$ 使 $\langle\gamma^n\rangle_{0,\beta}$ 非零，则 $\Phi\in\C[[y_0,\dots,y_m]]$，即 $y_0,\dots,y_m$ 的形式幂级数环。例如在 Fano 情形（即 $c_1(\cT_X)$ 丰沛）这成立，如式 (2.1) 所示。在非 Fano 情形，有时必须改在 Novikov 环上工作，但这里不这样做。因此在上述情形下，我们可以把 $\Phi$ 看作 $0\in H^{*}(X,\C)$ 的形式邻域上的函数。

当 $X$ 有非平凡的奇次上同调时，$\Phi$ 的定义有一些微妙之处，因为 $T_i\cup T_j=(-1)^{\deg T_i+\deg T_j}T_j\cup T_i$。这实质上是迫使我们把 $H^{*}(X,\C)$ 看作超流形。假设每个 $T_i$ 都有确定的次数，即 $T_i\in H^{d_i}(X,\C)$，且 $\deg T_i=d_i$。那么由 Gromov–Witten 不变量的定义，我们有
\[
\langle\dots,T_i,T_j,\dots\rangle_{g,\beta}=(-1)^{\deg T_i+\deg T_j}\langle\dots,T_j,T_i,\dots\rangle_{g,\beta}.
\]
若在 $\gamma=\sum_{i=0}^{m}y_iT_i$ 中把坐标 $y_i$ 取为\emph{超交换}的，即
\[
y_iy_j=(-1)^{\deg T_i+\deg T_j}y_jy_i,\qquad y_iT_j=(-1)^{\deg T_i+\deg T_j}T_jy_i,
\]
则 $y_iT_i$ 与 $y_jT_j$ 对一切 $i,j$ 交换。所以 $\Phi$ 应看作超交换变量的函数，且这一交换规则必须一致地应用。由于这里主要兴趣在 $X=\PP^2$（没有奇次上同调），我们常常限于 $X$ 只有偶次上同调的情形，以免过多担心符号。

利用 $\Phi$，我们定义 $X$ 的（大）量子上同调为环\footnote{在下一节中，我们将把量子上同调重新定义为某个流形切丛上的一族乘积。}
\[
H^{*}(X,\C[[y_0,\dots,y_m]]),
\]
% [Tutor Note: 原文 p.40 的符号约定 $y_iy_j=(-1)^{\deg T_i+\deg T_j}y_jy_i$ 只在 $X$ 有奇次上同调时才真正起作用；脚注所说"下一节重新定义量子上同调"即 §2.2（本块之后由 c02b 续译），那里把乘积看作某个流形切丛上的一族乘积。]

其乘积在生成元上由
\[
T_i*T_j=\sum_{k}(\partial_{y_i}\partial_{y_j}\partial_{y_k}\Phi)T^{k}
\]
给出，其中 $T^{0},\dots,T^{m}$ 是对偶于 $T_0,\dots,T_m$ 的 Poincaré 对偶基。常常有用的是定义 $g^{ij}$ 为矩阵
\[
g_{ij}=\int_X T_i\cup T_j
\]
的逆矩阵，
% [原文 p.41]
于是 $T^{i}=\sum g^{ij}T_j$，并且
\[
T_i*T_j=\sum_{a,k}(\partial_{y_i}\partial_{y_j}\partial_{y_a}\Phi)g^{ak}T_k.
\]

这个乘积是结合的，这一点并不显然；它结合是一个基本观察，可由 Gromov–Witten 不变量的附加性质证明。把等式
\[
(T_i*T_j)*T_k=T_i*(T_j*T_k)
\]
按上面的定义写开，就得到 WDVV 方程
\begin{align*}
&\sum_{a,b}(\partial_{y_i}\partial_{y_j}\partial_{y_a}\Phi)g^{ab}(\partial_{y_b}\partial_{y_k}\partial_{y_l}\Phi)\\
&\quad=(-1)^{\deg T_i(\deg T_j+\deg T_k)}\sum_{a,b}(\partial_{y_j}\partial_{y_k}\partial_{y_a}\Phi)g^{ab}(\partial_{y_b}\partial_{y_i}\partial_{y_l}\Phi),
\end{align*}
即一个偏微分方程组。

这里我们不证明结合性，因为许多量子上同调的基本文献已涵盖此事，例如 [28]。我们改为像 [28] 那样考查 WDVV 方程的一个标准推论，特别是关于 $\PP^2$ 的 Gromov–Witten 不变量。

先考虑 $\beta=0$ 时的 $\frac{1}{n!}\langle\gamma^n\rangle_{0,\beta}$。由点映射公理，除 $n=3$ 外它为零；而 $n=3$ 时
\[
\frac{1}{3!}\langle\gamma^3\rangle_{0,0}=\frac{1}{3!}\int_X\gamma\cup\gamma\cup\gamma,
\]
这纯粹是经典信息。于是把 $\Phi$ 分裂为
\[
\Phi=\Phi_{\text{classical}}+\Phi_{\text{quantum}},
\]
其中
\[
\Phi_{\text{quantum}}=\sum_{n\ge 0}\sum_{\substack{\beta\in H_2(X,\Z)\<本地路径> 0}}\frac{1}{n!}\langle\gamma^n\rangle_{0,\beta},
\]
而 $\Phi_{\text{classical}}$ 满足
\[
\partial_{y_i}\partial_{y_j}\partial_{y_k}\Phi_{\text{classical}}=\int_X T_i\cup T_j\cup T_k,
\]
从而对 $T_i*T_j$ 贡献
\[
\sum_{k}\left(\int_X T_i\cup T_j\cup T_k\right)T^{k}=T_i\cup T_j,
\]
即经典杯积。回忆我们取 $T_0$ 为 $X$ 的基本类、$T_1,\dots,T_p$ 为 $H^2(X,\C)$ 的生成元，于是对 $\beta\neq 0$ 可写
\[
\frac{1}{n!}\langle\gamma^n\rangle_{0,\beta}=\sum_{\substack{n_0,\dots,n_m\ge 0\\\sum n_j=n}}\epsilon(n_0,\dots,n_m)\langle T_0^{n_0},\dots,T_m^{n_m}\rangle_{0,\beta}\frac{y_0^{n_0}\cdots y_m^{n_m}}{n_0!\cdots n_m!},
\]
其中符号 $\epsilon(n_0,\dots,n_m)$ 由超交换关系
\[
(y_0T_0)^{n_0}\cdots(y_mT_m)^{n_m}=\epsilon(n_0,\dots,n_m)T_0^{n_0}\cdots T_m^{n_m}y_0^{n_0}\cdots y_m^{n_m}
\]
确定。

% [原文 p.42]
基本类公理告诉我们，上述和中只有 $n_0=0$ 时才有非零贡献；而除子公理给出
\begin{align}
&\sum_{n\ge 0}\frac{1}{n!}\langle\gamma^n\rangle_{0,\beta}\notag\\
&\quad=\sum_{n_1,\dots,n_m\ge 0}\epsilon(0,n_1,\dots,n_m)\left(\prod_{i=1}^{p}\frac{1}{n_i!}\left(y_i\int_{\beta}T_i\right)^{n_i}\right)\langle T_{p+1}^{n_{p+1}},\dots,T_m^{n_m}\rangle_{0,\beta}\cdot\notag\\
&\qquad\cdot\frac{y_{p+1}^{n_{p+1}}\cdots y_m^{n_m}}{n_{p+1}!\cdots n_m!}\notag\\
&\quad=\sum_{n_{p+1},\dots,n_m\ge 0}\epsilon(0,\dots,0,n_{p+1},\dots,n_m)\langle T_{p+1}^{n_{p+1}},\dots,T_m^{n_m}\rangle_{0,\beta}\notag\\
&\qquad\cdot\left(\prod_{i=1}^{p}e^{y_i\int_{\beta}T_i}\right)\frac{y_{p+1}^{n_{p+1}}\cdots y_m^{n_m}}{n_{p+1}!\cdots n_m!}. \tag{2.3}
\end{align}

% [Tutor Note: (2.3) 只作了一次编号（原文编号印在第二条等号那一行的左边空白处）；本译按全稿规则把编号置于多行显示的最后一行。第二个等号处把 $\sum_{n_i}\frac{1}{n_i!}(y_i\int_\beta T_i)^{n_i}$ 收成 $e^{y_i\int_\beta T_i}$，这正是 $\Phi_{\text{quantum}}$ 对 $H^2$ 变量取指数型的依据。]

我们对 $\PP^2$ 写出这一式子，其中 $T_0,T_1,T_2$ 分别是 $\PP^2$、一条直线与一个点的上同调类。我们有
\[
\Phi_{\text{quantum}}(y_0,y_1,y_2)=\sum_{d>0}\langle T_2^{3d-1}\rangle_{0,d[\ell]}e^{dy_1}\frac{y_2^{3d-1}}{(3d-1)!}.
\]
令 $\Gamma_{ijk}=\partial_{y_i}\partial_{y_j}\partial_{y_k}\Phi_{\text{quantum}}$，我们得到量子乘积的如下描述：
\begin{align*}
T_0*T_i&=T_i\\
T_1*T_1&=T_2+\Gamma_{111}T_1+\Gamma_{112}T_0\\
T_1*T_2&=\Gamma_{121}T_1+\Gamma_{122}T_0\\
T_2*T_2&=\Gamma_{221}T_1+\Gamma_{222}T_0
\end{align*}
其中
\begin{align*}
\Gamma_{111}&=\sum_{d\ge 1}d^3\langle T_2^{3d-1}\rangle_{0,d[\ell]}e^{dy_1}\frac{y_2^{3d-1}}{(3d-1)!}\\
\Gamma_{112}&=\sum_{d\ge 1}d^2\langle T_2^{3d-1}\rangle_{0,d[\ell]}e^{dy_1}\frac{y_2^{3d-2}}{(3d-2)!}\\
\Gamma_{122}&=\sum_{d\ge 1}d\langle T_2^{3d-1}\rangle_{0,d[\ell]}e^{dy_1}\frac{y_2^{3d-3}}{(3d-3)!}\\
\Gamma_{222}&=\sum_{d\ge 2}\langle T_2^{3d-1}\rangle_{0,d[\ell]}e^{dy_1}\frac{y_2^{3d-4}}{(3d-4)!}.
\end{align*}
记住 $\Phi_{\text{classical}}=\frac{y_0^2y_2}{2}+\frac{y_0y_1^2}{2}$，取 $i=j=1$、$k=l=2$ 的 WDVV 方程给出 $\Gamma_{222}+\Gamma_{111}\Gamma_{122}=\Gamma_{112}^2$，即
\[
\Gamma_{222}=\Gamma_{112}^2-\Gamma_{111}\Gamma_{122}.
\]

% [原文 p.43]
比较最后一个等式左、右两边 $y_2^{3d-4}$ 的系数，便得到 Vafa–Intriligator 公式（首次出现于 [60]）
\begin{align*}
&\langle T_2^{3d-1}\rangle_{0,d[\ell]}=\sum_{\substack{d_1,d_2>0\\d_1+d_2=d}}\langle T_2^{3d_1-1}\rangle_{0,d_1[\ell]}\langle T_2^{3d_2-1}\rangle_{0,d_2[\ell]}\cdot\\
&\qquad\cdot\left(d_1^2d_2^2\binom{3d-4}{3d_1-2}-d_1^3d_2\binom{3d-4}{3d_1-1}\right).
\end{align*}
由于 $\langle T_2^2\rangle_{0,[\ell]}=1$，其他一切亏格零 Gromov–Witten 不变量现在都由这个公式递推得到。所以 WDVV 方程在此情形下极为有力。

% [Tutor Note: Vafa–Intriligator 公式核对锚点：取 $d=2$，唯一取法 $d_1=d_2=1$，括号内为 $1\cdot1\binom{2}{1}-1\cdot1\binom{2}{0}=2-1=1$，故 $\langle T_2^5\rangle_{0,2[\ell]}=1$，与 p.38 Example 2.6 中"过 5 点唯一一条二次曲线"一致。括号中第二项为 $d_1^3d_2$（不是 $d_1^2d_2^2$），原式如此，勿"对称化"。]

\subsubsection{Frobenius 流形}

量子上同调环
\[
H^{*}(X,\C[[y_0,\dots,y_m]])
\]
可以看作一族环，由 $0\in H^{*}(X,\C)$ 的形式完备化参数化（若 $X$ 有奇次上同调，则是形式超流形）。注意到 $H^{*}(X,\C)$ 是 $H^{*}(X,\C)$ 在任一点处的切空间，我们可以把形式完备化看作一个形式流形 $\cM$，它带有 $T\cM$ 每个纤维上的环结构。这些环结构并非任意，而是具有使 $\cM$ 成为所谓 \emph{Frobenius 流形}的性质。

为避免关于符号约定的相当繁琐的讨论，从今往后我们假设 $X$ 没有奇次上同调，于是只需在通常流形（而非超流形）上定义 Frobenius 流形的概念。

以下 $\cM$ 将是一个复流形、复流形的芽，或一个非奇异概形。后面它还将是复流形沿某个复子流形的形式完备化。

\begin{definition}
设 $\cM$ 如上。$\cM$ 上的一个\emph{前 Frobenius 结构}是三元组数据 $(\nabla,g,\cA)$，其中
\begin{enumerate}[label=(\arabic*)]
\item $\nabla:\cT_\cM\to\cT_\cM\otimes\Omega^1_\cM$ 是平坦联络。
\item $g$ 是 $\cM$ 上的一个度量，即一个对称配对 $g:S^2(\cT_\cM)\to\cO_\cM$，它诱导同构 $\cT_\cM\cong\cT_\cM^{*}$。此外，$g$ 必须与 $\nabla$ 相容，即
\[
\dd(g(X,Y))=g(\nabla X,Y)+g(X,\nabla Y).
\]
\item $\cA:S^3(\cT_\cM)\to\cO_\cM$ 是一个对称张量。
\end{enumerate}
给定这些数据，就在 $\cM$ 的每个切空间上定义了乘积：$X\circ Y$ 由公式
\[
\cA(X,Y,Z)=g(X\circ Y,Z)
\]
定义。由 $\cA$ 的对称性，我们还有
\begin{equation}
g(X\circ Y,Z)=g(X,Y\circ Z). \tag{2.4}
\end{equation}
若这些数据还满足两条附加性质，则称这个前 Frobenius 流形为一个\emph{Frobenius 流形}。
\begin{enumerate}[label=(\arabic*),start=4]
\item 由 $\cA$ 定义的乘积是结合的。
\item 在 $\cM$ 上局部存在势函数 $\Phi$，使得
\[
\cA(X,Y,Z)=XYZ\Phi.
\]
\end{enumerate}
\end{definition}

% [原文 p.44]
\begin{example}
设 $X$ 是上同调只出现在偶次、且 $H^{*}(X,\C)$ 由 $T_0,\dots,T_m$ 生成的非奇异簇。若 $X$ 的 Gromov–Witten 势 $\Phi$ 属于 $\C[[y_0,\dots,y_m]]$（例如 $X$ 为 Fano 时），则 $\cM=\operatorname{Spec}\C[[y_0,\dots,y_m]]$ 给出一个 Frobenius 流形的例子。取 $\cM$ 上的常度量，由
\[
g(\partial_{y_i},\partial_{y_j})=\int_X T_i\cup T_j
\]
定义。联络 $\nabla$ 取平凡联络，以 $\partial_{y_0},\dots,\partial_{y_m}$ 为平坦截面。最后，
\[
\cA(\partial_{y_i},\partial_{y_j},\partial_{y_k})=\partial_{y_i}\partial_{y_j}\partial_{y_k}\Phi.
\]
此时乘积就是量子上同调乘积，即
\[
Y\circ Z=Y*Z.
\]
\end{example}

事实上，我们可以把量子上同调看作在某个大得多的流形上给出 Frobenius 流形结构。假设已选取 $H^2(X,\C)$ 的生成元 $T_1,\dots,T_p$，使它们实际上落在 $H^2(X,\R)$ 中，并且具有性质：对 $1\le i\le p$ 及任何可由稳定曲线表示的类 $\beta\in H_2(X,\Z)$，有 $\int_{\beta}T_i\ge 0$。例如，若 $H^{1,1}(X,\C)=H^2(X,\C)$，我们只需选取 $T_1,\dots,T_p$ 落在 $X$ 的 Kähler 锥的闭包中。然后我们引入新变量 $\kappa_1,\dots,\kappa_p$，满足关系
\[
e^{y_i}=\kappa_i,\qquad 1\le i\le p.
\]
现在假设对任何 $n_{p+1},\dots,n_m$，只有有限多个 $\beta$ 使得
\[
\langle T_{p+1}^{n_{p+1}},\dots,T_m^{n_m}\rangle_{0,\beta}\neq 0
\]
（例如 $X$ 为 Fano 时成立）。注意到 $\int_{\beta}T_1,\dots,\int_{\beta}T_p$ 全都非负，并有限多个选择地确定 $\beta$，由式 (2.3) 可见
\[
\Phi\in\C[\kappa_1,\dots,\kappa_p][[y_0,y_{p+1},\dots,y_m]].
\]
这个环是关于变量 $y_0,y_{p+1},\dots,y_m$、系数在多项式环 $\C[\kappa_1,\dots,\kappa_p]$ 中的形式幂级数环。

现在令
\begin{align}
\overline{\cM}&=\operatorname{Spec}\C[\kappa_1,\dots,\kappa_p][[y_0,y_{p+1},\dots,y_m]]\notag\\
\cM&=\operatorname{Spec}\C[\kappa_1^{\pm 1},\dots,\kappa_p^{\pm 1}][[y_0,y_{p+1},\dots,y_m]]. \tag{2.5}
\end{align}
先前使用的 $\cM=\operatorname{Spec}\C[[y_0,\dots,y_m]]$ 可以看作新 $\cM$ 中满足 $\kappa_1,\dots,\kappa_p=1$、$y_0=y_{p+1}=\cdots=y_m=0$ 的那一点的形式邻域。势 $\Phi$ 现在是 $\overline{\cM}$ 上的函数。注意 $\partial_{y_i}=\kappa_i\partial_{\kappa_i}$（$1\le i\le p$），于是我们先前定义的联络延拓为 $\cM$ 上以 $\partial_{y_0},\partial_{y_{p+1}},\dots,\partial_{y_m}$ 与 $\kappa_1\partial_{\kappa_1},\dots,\kappa_p\partial_{\kappa_p}$ 为平坦截面的平坦联络。这个联络不能跨过 $\overline{\cM}\setminus\cM$ 延拓。类似地，度量 $g$ 与对称形式 $\cA$ 也延拓到 $\cM$，其方式始终是用 $\partial_{y_i}=\kappa_i\partial_{\kappa_i}$（$1\le i\le p$）。

一般地，当我们在本章余下部分谈论量子上同调时，实际上指的是上述 $\cM$ 上或与之密切相关的流形上的 Frobenius 流形结构。

这个例子还有一些附加结构。

\begin{definition}
若 $\cM$ 是前 Frobenius 流形，则
\begin{enumerate}[label=(\arabic*)]
\item $\cM$ 上的向量场 $e$ 称为\emph{恒等}向量场，若对一切 $Y$ 有 $e\circ Y=Y$。
% [原文 p.45]
\item $\cM$ 上的向量场 $E$ 称为 \emph{Euler 向量场}，若对一切向量场 $Y$ 与 $Z$，
\begin{equation}
E(g(Y,Z))-g([E,Y],Z)-g(Y,[E,Z])=Dg(Y,Z) \tag{2.6}
\end{equation}
对某个常数 $D$ 成立，并且
\begin{equation}
[E,Y\circ Z]-[E,Y]\circ Z-Y\circ[E,Z]=d_0Y\circ Z \tag{2.7}
\end{equation}
对某个常数 $d_0$ 成立。
\end{enumerate}
\end{definition}

Euler 向量场可用来在向量场上产生分次：给定向量场 $Y$，若
\[
[E,Y]=dY,
\]
则称 $Y$ 是次数为 $d$ 的齐次向量场。

\begin{example}
沿用例 2.8，对生成元 $T_0,\dots,T_p$ 作同样的假设，于是我们在 $\cM=\operatorname{Spec}\C[\kappa_1^{\pm 1},\dots,\kappa_p^{\pm 1}][[y_0,y_{p+1},\dots,y_m]]$ 上有 Frobenius 流形结构。此时 $\partial_{y_0}$ 是平坦的恒等向量场：由 Gromov–Witten 不变量的基本类公理可知，$T_0$ 是量子上同调中的恒等元。

还存在一个 Euler 向量场，由
\[
E=\sum_{i=0}^{m}\left(1-\frac{\deg T_i}{2}\right)y_i\partial_{y_i}+\sum_{i=1}^{p}c_i\kappa_i\partial_{\kappa_i}
\]
定义，其中 $c_1(\cT_X)=\sum_{i=1}^{p}c_iT_i$。回忆对 $1\le i\le p$ 有 $\partial_{y_i}=\kappa_i\partial_{\kappa_i}$。

注意相对于这个向量场 $E$，形如
\[
Y=y_0^{a_0}\kappa_1^{a_1}\cdots\kappa_p^{a_p}y_{p+1}^{a_{p+1}}\cdots y_m^{a_m}\partial_{y_i}
\]
的向量场 $Y$ 是齐次的，其次数由以下方式给出：若 $i\notin\{1,\dots,p\}$，则 $y_i$ 的次数取为
\[
1-(\deg T_i)/2;
\]
若 $i\in\{1,\dots,p\}$，则 $\kappa_i$ 的次数为 $c_i$；而 $\partial_{y_i}$ 的次数为
\[
(\deg T_i)/2-1.
\]
这来自如下简单计算：
\[
[E,Y]=\left(-(1-(\deg T_i)/2)+\sum_{j=1}^{p}c_ja_j+\sum_{\substack{j=0<本地路径>,\dots,p\}}}^{m}(1-(\deg T_j)/2)a_j\right)Y.
\]
要验证 $E$ 是 Euler 向量场，注意到式 (2.6) 与式 (2.7) 的左、右两边都是张量，即都是 $\cO_\cM$-线性的，所以我们只需对 $Y=\partial_{y_i}$、$Z=\partial_{y_j}$ 验证。先验证式 (2.6)，我们看到
\begin{align*}
&E(g(\partial_{y_i},\partial_{y_j}))-g([E,\partial_{y_i}],\partial_{y_j})-g(\partial_{y_i},[E,\partial_{y_j}])\\
&\quad=g\left(\left(1-\frac{\deg T_i}{2}\right)\partial_{y_i},\partial_{y_j}\right)+g\left(\partial_{y_i},\left(1-\frac{\deg T_j}{2}\right)\partial_{y_j}\right)\\
&\quad=\left(2-\frac{\deg T_i+\deg T_j}{2}\right)g(\partial_{y_i},\partial_{y_j})\\
&\quad=(2-\dim_{\C}X)g(\partial_{y_i},\partial_{y_j}),
\end{align*}
因为仅当 $\deg T_i+\deg T_j=2\dim_{\C}X$ 时 $g(\partial_{y_i},\partial_{y_j})\neq 0$。于是 $D=2-\dim_{\C}X$。

% [Tutor Note: 上面 $[E,Y]$ 公式中第一项是 $-(1-(\deg T_i)/2)$（对 $Y$ 的 $\partial_{y_i}$ 部分求导所致），与各项 $a_j$ 的符号约定配套；若把第一项写成 $+$，会与最终"$\partial_{y_i}$ 次数为 $(\deg T_i)/2-1$"矛盾，可据此校验。]

% [原文 p.46]
下面我们断言 (2.7) 对 $d_0=1$ 成立。仍取 $Y=\partial_{y_i}$、$Z=\partial_{y_j}$，注意 $Y$ 与 $Z$ 分别是次数 $(\deg T_i)/2-1$ 与 $(\deg T_j)/2-1$ 的齐次向量场。于是为使式 (2.7) 成立，$Y\circ Z$ 必须是次数 $1+((\deg T_i)/2-1)+((\deg T_j)/2-1)$ 的齐次向量场。为看出这一点，先注意 $\Phi$ 本身是齐次的。事实上，由式 (2.3)，$\Phi$ 中的一项形如
\[
\langle T_{p+1}^{n_{p+1}},\dots,T_m^{n_m}\rangle_{0,\beta}\left(\prod_{i=1}^{p}\kappa_i^{\int_{\beta}T_i}\right)\frac{y_{p+1}^{n_{p+1}}\cdots y_m^{n_m}}{n_{p+1}!\cdots n_m!},
\]
其次数为
\begin{align*}
&\sum c_i\int_{\beta}T_i+\sum_{j=p+1}^{n}\left(1-\frac{\deg T_j}{2}\right)n_j\\
&\quad=\int_{\beta}c_1(\cT_X)+\sum_{j=p+1}^{n}n_j-\sum_{j=p+1}^{n}\left(\frac{\deg T_j}{2}\right)n_j.
\end{align*}
另一方面，由式 (2.1)，该项中的 Gromov–Witten 不变量为零，除非
\[
\int_{\beta}c_1(\cT_X)+\sum_{j=p+1}^{n}n_j-\sum_{j=p+1}^{n}\left(\frac{\deg T_j}{2}\right)n_j=3-\dim_{\C}X.
\]
因此 $\Phi$ 可以看作次数为 $3-\dim_{\C}X$ 的齐次量，而 $\partial_{y_i}\partial_{y_j}\partial_{y_k}\Phi$ 的次数为
\[
3-\dim_{\C}X+((\deg T_i)/2-1)+((\deg T_j)/2-1)+((\deg T_k)/2-1).
\]
综合以上，
\[
\partial_{y_i}\circ\partial_{y_j}=\sum_{k,l}(\partial_{y_i}\partial_{y_j}\partial_{y_k}\Phi)g^{kl}\partial_{y_l}
\]
的次数为
\begin{align*}
&3-\dim_{\C}X+((\deg T_i)/2-1)+((\deg T_j)/2-1)\\
&\qquad+((\deg T_k)/2-1)+((\deg T_l)/2-1)\\
&\quad=3-\dim_{\C}X-((\deg T_i)/2-1)+((\deg T_j)/2-1)\\
&\qquad-2+\dim_{\C}X\\
&\quad=1+((\deg T_i)/2-1)+((\deg T_j)/2-1),
\end{align*}
这是因为仅当 $\deg T_k+\deg T_l=2\dim_{\C}X$ 时 $g^{kl}\neq 0$。这就证明了式 (2.7)。\hfill$\square$
\end{example}

% [Tutor Note: 原文 p.46 此处的中间一行印作 "= 3 − dim_C X − ((deg T_i)/2 − 1) + ((deg T_j)/2 − 1)"，但由上一行到下一行（结果为 $1+((\deg T_i)/2-1)+((\deg T_j)/2-1)$）可反推出该行第一处符号应为 "+"，即 $\deg T_k+\deg T_l=2\dim_{\C}X$ 代入后应为 $3-\dim_{\C}X+((\deg T_i)/2-1)+((\deg T_j)/2-1)-2+\dim_{\C}X$。本译照原文印刷符号抄录（保留 "−"），仅在此注明正确形式，正文不作改动。]

若 $\cM$ 是带向量场 $E$ 且 $d_0\neq 0$ 的前 Frobenius 流形，我们可以在 $\cM\times\C^{\times}$ 上的向量丛 $p_1^{*}\cT_\cM$ 上定义联络 $\widehat\nabla$，其中 $p_1:\cM\times\C^{\times}\to\cM$ 是第一投影，$\C^{\times}=\C\setminus\{0\}$。设 $\hbar$ 是 $\C^{\times}$ 上的坐标。这个联络由
\begin{align*}
\widehat\nabla_X(Y)&=\nabla_X(Y)+\hbar^{-1}X\circ Y\\
d_0\widehat\nabla_{\hbar\partial_{\hbar}}(Y)&=\hbar\partial_{\hbar}Y-\hbar^{-1}E\circ Y+\operatorname{Gr}_E(Y)
\end{align*}
定义，其中 $\operatorname{Gr}_E$ 是 $\cO_\cM$-线性映射，它在平坦向量场 $Y$ 上由 $Y\mapsto[E,Y]$ 给出。这个联络称为\emph{第一结构联络}，或 \emph{Dubrovin 联络}。它有下列重要性质：

% [Tutor Note: Dubrovin 联络公式中 $\hbar$ 是 $\C^{\times}$ 上的坐标；原文在 p.46 末尾只给出定义与"具有下列重要性质："一句，续文在 p.47（下一块 c02b 开头），本块照原文止于此，未加内容。]
